% !TEX TS-program = pdflatex
% !TEX encoding = UTF-8 Unicode

% This is a simple template for a LaTeX document using the "article" class.
% See "book", "report", "letter" for other types of document.

\documentclass[11pt]{article} % use larger type; default would be 10pt

\usepackage[utf8]{inputenc} % set input encoding (not needed with XeLaTeX)

%%% Examples of Article customizations
% These packages are optional, depending whether you want the features they provide.
% See the LaTeX Companion or other references for full information.

%%% PAGE DIMENSIONS
\usepackage[margin=1in]{geometry} % to change the page dimensions
\geometry{a4paper} % or letterpaper (US) or a5paper or....
% \geometry{margin=2in} % for example, change the margins to 2 inches all round
% \geometry{landscape} % set up the page for landscape
%   read geometry.pdf for detailed page layout information

\usepackage{graphicx} % support the \includegraphics command and options

% \usepackage[parfill]{parskip} % Activate to begin paragraphs with an empty line rather than an indent

%%% PACKAGES
\usepackage{booktabs} % for much better looking tables
\usepackage{array} % for better arrays (eg matrices) in maths
\usepackage{paralist} % very flexible & customisable lists (eg. enumerate/itemize, etc.)
\usepackage{verbatim} % adds environment for commenting out blocks of text & for better verbatim
\usepackage{subfig} % make it possible to include more than one captioned figure/table in a single float
% These packages are all incorporated in the memoir class to one degree or another...

%%% HEADERS & FOOTERS
\usepackage{fancyhdr} % This should be set AFTER setting up the page geometry
\pagestyle{fancy} % options: empty , plain , fancy
\renewcommand{\headrulewidth}{0pt} % customise the layout...
\lhead{}\chead{}\rhead{}
\lfoot{}\cfoot{\thepage}\rfoot{}

%%% SECTION TITLE APPEARANCE
\usepackage{sectsty}
\allsectionsfont{\sffamily\mdseries\upshape} % (See the fntguide.pdf for font help)
% (This matches ConTeXt defaults)

%%% ToC (table of contents) APPEARANCE
\usepackage[nottoc,notlof,notlot]{tocbibind} % Put the bibliography in the ToC
\usepackage[titles,subfigure]{tocloft} % Alter the style of the Table of Contents
\renewcommand{\cftsecfont}{\rmfamily\mdseries\upshape}
\renewcommand{\cftsecpagefont}{\rmfamily\mdseries\upshape} % No bold!

\usepackage[backend=biber,style=chem-acs]{biblatex}
\addbibresource{srprefs.bib}
%\bibliographystyle{plain} % We choose the "plain" reference style
%\bibliography{refs} % Entries are in the refs.bib file

% For \onehalfspacing and \singlespacing macros
\usepackage{setspace}

\usepackage{etoolbox}

\AtBeginEnvironment{quote}{\par\singlespacing\small} % block quote

\usepackage{amsmath}
\usepackage{braket}

\usepackage{hyperref}
\hypersetup{colorlinks=true, citecolor=blue, urlcolor=blue, linkcolor=blue}
\usepackage{cleveref}	
  	\crefname{figure}{Figure}{Figures}
  	\crefname{table}{Table}{Tables}
  	\crefname{equation}{Eq.}{Eqs.}
  	\crefname{section}{Section}{Sections}
  	\crefname{subsection}{Section}{Sections}
  	\crefname{subsubsection}{Section}{Sections}
  	\crefname{algorithm}{Algorithm}{Algorithms}

%%% END Article customizations

%%% The "real" document content comes below...

\title{Common Quandry on the Justice of God}
\author{Samuel Powell}
\date{2025-02-11} % Activate to display a given date or no date (if empty),
         % otherwise the current date is printed 

\begin{document}
\maketitle
```
If a child-molesting person is on his deathbed and the confesses and repents and gives his life to Christ before he dies he goes to heaven right? What about the 12 year old boy that he raped, what if he doesn't believe in God, does that mean he goes to hell while the child molester goes to heaven?
```

Technical answer: yes. The real question is usually in the form of an objection to or questioning this. Some dangers are equating suffering with moral capital (suffering is not (evidently) distributed evenly, but all do and none deserve or have a right to expect better). The forgiveness of the repentant is `yes, and thank God for it'. We also are tempted to read `doesn't believe in God' as `wasn't properly convinced' or `just didn't see the issue clearly'. These are challenging, but the real issue is `chose not to submit to the authority of God' or `did not love God and choose to pursue relationship or union with Him'. The ultimate issue, I think (in salvation), is will and desire, not belief or doctrine. Believe and doctrine is something you can `just get wrong' and kinda depends on what's presented to you, the orientation of your will and your heart is yours alone. Think also: some will be clearly presented with the authority and nature of God and choose to reject Him and the life He as for us. (Revelation 20, fall of satan and angels, maybe the rich man in Luke 16:19-31).

There are two obvious, major paths of objection to this: 1) why would God forgive and accept the rapist? 2) why wouldn't God forgive and accept the child? There is also just the objection to the contrast between the two, but this seems like more of an emotional or subjective reaction, not a reasoned issue. Not to say it's negligible, but it can probably be reduced to one or both of the aforemetioned routes on consideration and meditation. 

For the first: God is gracious and merciful, there is no record of one who would repent, yet is rejected by God. Take the two `robbers' on the crosses beside Jesus: (Luke 23:39-43) there is no difference indicated in their records or the terms of their condemnation, only in their response to Christ: one self-seeking, only hoping to use Christ to escape his fate, the other seeking to be a part of Jesus' kingdom, submitting to His authority. Consider also 2 Samuel 11, Genesis 38, Genesis 34, the apostle Paul. Not all of these are necessarily `Christian' or `saved' or in good relationship with God, but God certainly used them all. David was. 

On the second objection, it's important to remember that rebellion against God does not necessitate obviously evil behavior, yet it is that rebellion that one is ultimately judged on / by. Consider Apollos in Acts 18:24-28, not well versed in theology, but fervently teaching (this example isn't relevant to what I'm saying here). One of the strongest or most sobering examples is in Matthew 7:15-23 (Luke 6:46, Num. 24:4). These even did major `spiritual' works \emph{for Christ} (at least nominally) but didn't know him, and are rejected. Those who reject the authority of Christ are no less likely to suffer at the hand / according to the ravages of sin, but they reject the free offer of joining His kingdom. Citizenship is not earned, not by works of righteousness or by warrant of suffering. 


\printbibliography

\end{document}
